% Generated draft: requires mathematical and source review.
\documentclass[11pt]{article}
\usepackage[margin=25mm]{geometry}
\usepackage{fontspec}
\setmainfont{Latin Modern Roman}
\usepackage{amsmath,amssymb,mathtools,unicode-math}
\setmathfont{Latin Modern Math}
\usepackage{enumitem,needspace,xurl,hyperref}
\setlength{\parindent}{0pt}
\setlength{\parskip}{5pt}
\setlength{\emergencystretch}{3em}
\newcommand{\N}{\mathbb{N}}
\newcommand{\Z}{\mathbb{Z}}
\newcommand{\Q}{\mathbb{Q}}
\newcommand{\R}{\mathbb{R}}
\newcommand{\C}{\mathbb{C}}
\newcommand{\F}{\mathbb{F}}
\newcommand{\A}{\mathbb{A}}
\newcommand{\PP}{\mathbb{P}}
\newcommand{\E}{\mathbb{E}}
\newcommand{\eps}{\varepsilon}
\newcommand{\dd}{\,\mathrm{d}}
\begin{document}
\section*{UnsolvedMath ID: 10800012}
\textbf{Source problem code:} catalogue entry
\subsection{Definitions and mathematical statement}
For a smooth compact variety $X\subset\R^N$ without boundary, a conventional interpretation of the stated complexity is
\[
 c_{\mathrm{lin}}(X)=
 \min_{\substack{a\in(\R^N)^*\\ a|_X\ \mathrm{Morse}}}
 \#\operatorname{Crit}(a|_X).
\]
Here a Morse function has only nondegenerate critical points. This formula is conditional on the displayed regularity and compactness assumptions; the excerpt does not actually specify the class of varieties, and singular or noncompact cases require additional conventions. The embedding in $\R^N$ is part of the data: this is not the minimum over all smooth Morse functions on the abstract underlying manifold.

Vassiliev's Problem 5B asks for the version of the preceding simplification problem using this linear-function critical-point complexity. The permitted surgeries affecting the measure should include not only topological changes of the variety but also bifurcations of its dual variety. In the usual geometric interpretation, the dual records tangent hyperplanes, so its behavior is relevant to critical points of restricted linear functions.

The preceding problem is not supplied. Consequently the required simplification criterion, the class of admissible varieties and deformations, and the exact allowed surgeries cannot be reconstructed from this record. The registry calls this a greedy-simplification question, but that phrase alone does not specify a theorem about the existence, termination, or monotonicity of a simplification procedure.
\subsection{Short English statement}
Study the preceding simplification problem using the minimum number of critical points among Morse restrictions of ambient linear functions, allowing relevant dual-variety bifurcations as well as topological surgeries.
\subsection{Sources}
\begin{itemize}
\item Local registry: \url{data/unsolvedmath/problems.json}, permanent UnsolvedMath ID $10800012$, exact source code \texttt{catalogue entry}, attributed to Victor Vassiliev.
\item Vassiliev, \emph{A Few Problems on Monodromy and Discriminants}, Problem 5B: \url{https://amj.math.stonybrook.edu/html-articles/Files-2015-2024/15-11/}. Supplied access date: 2026-07-29; extraction: html-structure; eighteen labeled problem environments.
\item The excerpt explicitly depends on the previous problem, which is absent from this record. No additional primary reference is supplied for the registry's claim that no solution or partial result was found.
\end{itemize}
\end{document}
